% NOTE: this document uses the Springer LNCS class (llncs.cls) and the
% splncs04 bibliography style. On Overleaf, start from the official
% "Lecture Notes in Computer Science" template (or upload llncs.cls /
% splncs04.bst alongside this file) so it compiles without missing-class
% errors.
\documentclass[runningheads]{llncs}

% ---- Packages ----
\usepackage[T1]{fontenc}
\usepackage{amsmath, amssymb}
\usepackage{graphicx}
\usepackage{booktabs}
\usepackage{multirow}
\usepackage{subcaption}
\usepackage{xcolor}
\usepackage{hyperref}
\hypersetup{colorlinks=true, linkcolor=black, citecolor=black, urlcolor=blue}
\renewcommand{\arraystretch}{0.92}
\setlength{\textfloatsep}{8pt plus 2pt minus 2pt}
\setlength{\floatsep}{6pt plus 2pt minus 2pt}
\setlength{\intextsep}{6pt plus 2pt minus 2pt}
\setlength{\abovecaptionskip}{4pt}
\setlength{\belowcaptionskip}{2pt}

\newcommand{\TODO}[1]{\textcolor{red}{\textbf{[TODO: #1]}}}

\begin{document}

\title{Multi-View Basketball Player and Ball Tracking with 3D Reconstruction}

\titlerunning{Multi-View Player and Ball Tracking with 3D Reconstruction}

\author{Andrea Renelli \and Marco Misseroni}

\authorrunning{A. Renelli and M. Misseroni}

\institute{University of Trento, Trento, Italy\\
\email{andrea.renelli@studenti.unitn.it}, \email{marco.misseroni@studenti.unitn.it}}

\maketitle

\begin{abstract}
We present a multi-camera pipeline for tracking players, referees and the
ball during a basketball action recorded at the Sanbapolis facility, and
for reconstructing their 3D trajectories via triangulation. Ground truth
is annotated at 5\,fps on three synchronized views ($2846$ boxes, $309$
frames, $16$ categories). Players/referees are tracked per camera at the
native 25\,fps with a YOLO detector and its built-in tracker, reaching
Precision~$=0.775$, Recall~$=0.712$ (F1~$=0.742$), MOTP~$=0.423$ across
all three views. The ball is far harder: a classical motion/colour
baseline and a zero-shot pretrained small-ball detector (WASB/HRNet, with
tiling on the widest-shot camera) reach F1~$=0.068$ and F1~$=0.143$
respectively -- the latter roughly double, but both still limited by the
ball's small size (a first fine-tuned-detector attempt failed for the same
structural reason and was abandoned). Camera intrinsics/extrinsics were
provided for this project; 2D tracks are re-identified across views by a
reprojection-error voting scheme and triangulated into a court-centered
frame. Against pseudo-ground truth from triangulated annotations, the
reconstruction reaches a median error of $0.72$\,m (RMSE $1.01$\,m) over
$44.9\%$ of candidate points, and the 18 resulting 3D tracks keep a valid
position in $88.1\%$ of the frames within their own span on average.
\end{abstract}

\section{Introduction}
\label{sec:intro}

This work addresses the Player Tracking project of the Computer Vision
course: given a multi-camera setup at the Sanbapolis facility, detect and
track players, referees and the ball in each view, and reconstruct their
3D trajectories by triangulation. One action was recorded from three
fixed, synchronized cameras (\texttt{out2}, \texttt{out4}, \texttt{out13})
on a FIBA-regulation $28\times15$\,m court. As a team of two, the
mandatory scope is annotation (task~1), 2D tracking with evaluation
(task~2/2a), and 3D triangulation with evaluation (task~3/3a) using at
least two views; we use all three throughout.

Two aspects make this work not straight-forward. The ball is
small and fast, frequently occluded or blurred, and measures well under
$20$\,px after the mandatory resize to a detector's input resolution --
comparable to a CNN backbone's stride. And cross-camera association
cannot reuse a per-camera ID: correspondence has to come from projective
geometry, since no appearance-based re-identification was used (future
work, Sec.~\ref{sec:conclusion}).

This report covers the dataset, the camera model, the 2D
detection/tracking pipeline (players and ball), and cross-camera 3D
triangulation, each evaluated quantitatively -- including a
trajectory-level continuity measure -- against the annotated ground
truth (Sec.~\ref{sec:results}).

\section{Method}
\label{sec:method}

\subsection{Dataset and Annotation}
\label{sec:dataset}

One action was annotated from the three views on Roboflow at 5\,fps:
$309$ frames ($99$/$105$/$105$ for \texttt{out2}/\texttt{out4}/
\texttt{out13}) and $2846$ boxes over $16$ categories -- a generic
\texttt{person} catch-all, \texttt{Ball} ($191$ boxes), six jersey-numbered
``Red'' identities and six ``White'' (two of the twelve,
\texttt{Red\_13}/\texttt{White\_11}, appear in only $1$/$5$ frames and are
edge cases rather than consistently tracked identities), and two referees.
Per-identity annotation lets the 3D pseudo-ground-truth
(Sec.~\ref{sec:3d_results}) match a specific person across cameras without
a separate re-identification step; for the 2D evaluation
(Sec.~\ref{sec:2d_results}) every non-ball category is pooled into one
\texttt{player} class.

\subsection{Camera Model}
\label{sec:calibration}

Per-camera intrinsics ($K$, distortion) and extrinsics (rotation/
translation w.r.t.\ a court-centered frame) were provided for this
project rather than self-calibrated. They undistort every frame before 2D
detection (Sec.~\ref{sec:2d_tracking}) and build each camera's projection
matrix $P=K[R\,|\,t]$ for triangulation (Sec.~\ref{sec:3d_tracking}); we
do not report a separate reprojection-error figure for the calibration
itself, since it was not computed by us -- the full undistortion+
triangulation chain is instead validated end-to-end in
Sec.~\ref{sec:3d_results}.

\subsection{2D Detection and Tracking}
\label{sec:2d_tracking}

\paragraph{Players and referees.} Each camera is processed independently
at the native $25$\,fps: frames are undistorted (a per-pixel remap built
once, reused every frame), resized to $1920\times1088$, and passed to a
YOLO detector (COCO \texttt{person}/\texttt{sports ball}
classes~\cite{Jocher2023}) run through Ultralytics' tracking mode
(\texttt{persist=True}), which defaults to BoT-SORT~\cite{Aharon2022} when
no alternative tracker is configured, as here. A low confidence threshold
($0.1$) favors recall, relying on the tracker's temporal persistence to
filter spurious detections.

\paragraph{Ball.} The same pipeline is largely blind to the ball: at
$1920\times1088$ it covers a handful of pixels, below what the backbone
preserves after downsampling. A first attempt fine-tuned YOLO on ball-only
crops with SAHI tiled inference~\cite{Akyon2022}; this did not fix recall,
which is structural (the ball shrinks to single-digit pixel counts even
per tile) rather than a training issue. Two alternatives were evaluated
instead: (a) a classical baseline (MOG2 background subtraction with
shadow removal, morphological cleanup, a circularity/aspect-ratio shape
filter and an HSV colour-pattern filter, per-camera area thresholds); and
(b) WASB/HRNet~\cite{Tarashima2023}, a pretrained heatmap-based detector
with a basketball-specific checkpoint, used zero-shot with a confidence
threshold recalibrated empirically (well below the broadcast-tuned
default). On the widest-shot camera (\texttt{out13}, ball $\approx\!7$\,px
after resize) WASB is additionally run with $3\times3$ overlapping tiling.

\subsection{Cross-Camera Association and 3D Triangulation}
\label{sec:3d_tracking}

Per-camera tracks are undistorted using the calibration from
Sec.~\ref{sec:calibration}. For every camera pair and frame, candidate
matches between per-camera track IDs are scored by triangulating each
detection pair and measuring the combined reprojection error in both
views~\cite{Hartley2004}; a Hungarian assignment selects the best pairing
per frame (reprojection error $<30$\,px). A track-level global identity
forms only from pairs that are \emph{mutually} the best match and win the
per-frame assignment in $\geq5$ distinct frames; consistent pairs merge
into cross-camera groups via union-find, one global ID per group.

For every frame, detections sharing a global ID from $\geq2$ cameras are
triangulated with a multi-view DLT solve. A court-sized plausibility box
($28\times15$\,m $+5$\,m margin, height $[-1,3]$\,m) discards points
outside it, removing residual wrong matches rather than genuine players.

\section{Results}
\label{sec:results}

\subsection{Experimental Setup}
\label{sec:exp_setup}

The 2D evaluation matches predicted/ground-truth boxes per frame (per
camera, per class) with a Hungarian assignment on IoU, reporting
Precision, Recall, F1 and MOTP (mean IoU over matched pairs), following
CLEAR MOT~\cite{Bernardin2008}. The acceptance threshold is deliberately
low (any non-zero overlap counts as a candidate match), so
Precision/Recall/F1 mainly reflect whether an object was found in roughly
the right place, while MOTP carries the localization-quality signal. The
3D evaluation triangulates annotated boxes ($\geq2$ views, same frame,
same category) into a pseudo-ground-truth 3D point, matched to the
reconstruction with a Hungarian assignment on 3D distance (threshold
$3$\,m).

\subsection{2D Tracking}
\label{sec:2d_results}

Table~\ref{tab:2d_players} reports player/referee accuracy per camera and
the ball trackers pooled over all three (per-camera ball numbers follow
the same \texttt{out13}-is-hardest pattern). \texttt{out2}/\texttt{out4}
(closer, narrower shots) track players considerably better than
\texttt{out13} (wide, more/smaller/more-occluded people); pooled,
Recall~$=0.712$, Precision~$=0.775$, MOTP~$=0.423$. For the ball, WASB
roughly doubles F1 over the classical baseline with fewer false
positives, but both miss most instances -- even WASB finds fewer than 1
in 6 -- consistent with the small-object argument above.

\begin{table}[t]
\centering\small
\caption{2D tracking accuracy: players+referees pooled (class\_id 0) per
camera and aggregated; ball trackers (class\_id 1) pooled over all three
cameras.}
\label{tab:2d_players}
\begin{tabular}{llccccc}
\toprule
& & TP & FP & FN & Precision / Recall / F1 & MOTP \\
\midrule
\multirow{4}{*}{Players}
& \texttt{out2}  & 583  & 44  & 179 & 0.930 / 0.765 / 0.839 & 0.441 \\
& \texttt{out4}  & 459  & 98  & 172 & 0.824 / 0.727 / 0.773 & 0.470 \\
& \texttt{out13} & 848  & 408 & 414 & 0.675 / 0.672 / 0.674 & 0.386 \\
& \textbf{All}   & 1890 & 550 & 765 & \textbf{0.775 / 0.712 / 0.742} & \textbf{0.423} \\
\midrule
\multirow{2}{*}{Ball}
& Classic (MOG2+HSV) & 16 & 263 & 175 & 0.057 / 0.084 / 0.068 & 0.116 \\
& WASB/HRNet         & 25 & 133 & 166 & 0.158 / 0.131 / 0.143 & 0.083 \\
\bottomrule
\end{tabular}
\end{table}

\subsection{3D Reconstruction}
\label{sec:3d_results}

Of $1085$ pseudo-ground-truth 3D points, $487$ ($44.9\%$ coverage) match a
reconstructed point within $3$\,m, with an overall MED of $0.722$\,m and
RMSE of $1.013$\,m. Per-identity (Table~\ref{tab:3d_breakdown}), most
players/referees localize with sub-meter accuracy, with two exceptions:
the ball (only $3$ matches, limited by the same detection scarcity as
Sec.~\ref{sec:2d_results}) and \texttt{Red\_23} ($1.91$\,m over $75$
matches -- not a handful of outliers, pointing to a systematic
cross-camera association problem for that identity).

\begin{figure}[h]
\centering
\includegraphics[width=0.42\linewidth]{3d_positions.png}
\caption{Reconstructed 3D trajectories, court-centered frame; each colour
is one global cross-camera identity (Sec.~\ref{sec:3d_tracking}).}
\label{fig:traj3d}
\end{figure}

\begin{table}[h]
\centering\small
\caption{3D error by identity, sorted by mean error (two columns for
space).}
\label{tab:3d_breakdown}
\begin{tabular}{lcc@{\hspace{1.2cm}}lcc}
\toprule
Identity & Mean [m] & $n$ & Identity & Mean [m] & $n$ \\
\midrule
\texttt{Red\_2}    & 0.307 & 29 & \texttt{White\_14} & 0.592 & 26 \\
\texttt{Refree\_1} & 0.308 & 49 & \texttt{White\_2}   & 0.614 & 32 \\
\texttt{Red\_9}    & 0.381 & 14 & \texttt{Red\_7}     & 0.668 & 11 \\
\texttt{White\_27} & 0.460 & 45 & \texttt{White\_34}  & 0.678 & 71 \\
\texttt{White\_16} & 0.476 & 82 & \texttt{Ball}       & 1.431 & 3  \\
\texttt{Red\_11}   & 0.496 & 50 & \texttt{Red\_23}    & 1.907 & 75 \\
\bottomrule
\end{tabular}
\end{table}

\subsection{Trajectory Continuity}
\label{sec:traj_continuity}

Point accuracy says nothing about identity continuity over time, which
task~3a asks for explicitly. For each of the $18$ global 3D tracks,
\emph{continuity} is the fraction of frames within its own
first-to-last-frame span with a reconstructed position ($1.0$ = no
gaps): $88.1\%$ on average (median span $155$ frames $=6.2$\,s; mean
$3.4$ gaps/track). More tellingly, $18$ IDs exceeds the $\sim\!12$--$14$
people on court (10 players + 2 referees): vote-based re-identification
(Sec.~\ref{sec:3d_tracking}) sometimes fragments one person into several
global tracks once cross-camera correspondence is lost, and never
re-merges them -- a limitation of an association scheme with no
appearance model.

\section{Conclusions and Future Work}
\label{sec:conclusion}

We built and evaluated a three-camera pipeline for player/referee and
ball tracking with 3D reconstruction, covering the mandatory tasks for a
two-person team using all three views rather than the two required.
Players/referees reach moderate-to-good 2D accuracy (pooled F1~$=0.742$)
and sub-meter 3D median error ($0.72$\,m), validating the pipeline for
the common case; the ball and \texttt{Red\_23} are the two clear
outliers. The ball proved hard throughout: a fine-tuned-detector-plus-
tiling attempt was abandoned once diagnosed as structurally limited by
object size, and even the better pretrained model (WASB) stays well
short of usable.

Future work: diagnose the \texttt{Red\_23} outlier (likely a recurring
wrong cross-camera pairing); add appearance-based re-identification
against the fragmentation of Sec.~\ref{sec:traj_continuity}; revisit ball
detection at its true pixel scale; adopt identity-aware MOT metrics
(MOTA, ID switches); and, as a bonus, import the motion into Unreal
Engine (task~4).

\small
\begin{thebibliography}{9}

\bibitem{Hartley2004}
R.~Hartley and A.~Zisserman,
\emph{Multiple View Geometry in Computer Vision}, 2nd ed.
Cambridge University Press, 2004.

\bibitem{Akyon2022}
F.~C.~Akyon, S.~O.~Altinuc, and A.~Temizel,
``Slicing Aided Hyper Inference and Fine-tuning for Small Object
Detection,''
in \emph{Proc. IEEE Int. Conf. on Image Processing (ICIP)}, 2022,
pp.\ 966--970.

\bibitem{Aharon2022}
N.~Aharon, R.~Orfaig, and B.-Z.~Bobrovsky,
``BoT-SORT: Robust Associations Multi-Pedestrian Tracking,''
arXiv preprint arXiv:2206.14651, 2022.

\bibitem{Tarashima2023}
S.~Tarashima, M.~A.~Haq, Y.~Wang, and N.~Tagawa,
``Widely Applicable Strong Baseline for Sports Ball Detection and
Tracking,''
in \emph{Proc. British Machine Vision Conference (BMVC)}, 2023.

\bibitem{Bernardin2008}
K.~Bernardin and R.~Stiefelhagen,
``Evaluating Multiple Object Tracking Performance: The CLEAR MOT
Metrics,''
\emph{EURASIP J. Image and Video Processing}, vol.\ 2008, 2008.

\bibitem{Jocher2023}
G.~Jocher, A.~Chaurasia, and J.~Qiu,
``Ultralytics YOLO,'' 2023. [Software]
\url{https://github.com/ultralytics/ultralytics}

\end{thebibliography}

\end{document}
